\chapter{Magnetism}
\setwlchapter{magnetism}
\section{Exchange on a cluster}
A positive $J$ in $J\vS_i\cdot\vS_j$ favours a singlet. Superexchange
from a half-filled one-band model has this sign, whereas a direct
Hund interaction favours alignment. In double exchange a mobile
carrier aligned with local moments instead gains kinetic energy
when neighbouring moments are parallel. These mechanisms need not
produce the same dependence on filling or on bond angle.

For three spins $1/2$ on an equilateral triangle all pair bonds are
present, so the energy depends only on the total spin:
$H=J[\mathbf S_\mathrm{tot}^2-9/4]/2$. Its two low-spin doublets
have the same energy even though no state is a singlet on all three
bonds at once. On a square ring not all pairs are neighbours;
the Casimir alone no longer fixes each energy.
\casestudy{frustration}{Resolving the four-site ring by magnetization}

\section{Spin waves and their regime of validity}
Expanding around a ferromagnet polarized along $+z$ gives
$S^z=s-b^\dagger b$, $S^+\simeq\sqrt{2s}\,b$. In
$-J\sum_i\vS_i\cdot\vS_{i+1}$ the quadratic magnon energy is
$2Js(1-\cos k)$; the absent gap follows from the global spin-rotation
symmetry. For a Néel reference state rotate one sublattice first.
Opposite classical spins then become parallel local axes, and the
transverse exchange generates \emph{pairs} of bosons. Diagonalizing
the bosonic quadratic form requires the commutation metric with
opposite signs on particles and holes; a unitary diagonalization of
the Hermitian energy matrix does not give the magnon frequencies.
\casestudy{magnons}{A staggered field as a gap regulator}
The quadratic truncation assumes a small mean occupation relative to
$2s$. A gapless antiferromagnetic chain of spin $1/2$ is particularly
unfavourable for a broken-symmetry expansion; finite positive field
in the example is part of the Hamiltonian, not just a numerical tolerance.

\section{One dimension: strings and boundaries}
The Jordan--Wigner string counts the occupied sites to the left of
each spin flip. Adjacent strings cancel in the bulk. Long-range flips
retain intervening density factors, and closing a chain into a ring
introduces a boundary sign that depends on total fermion parity.
The open XXZ chain is therefore a useful spectrum-level test before
handling periodic boundary conditions.
\casestudy{xxz}{The interacting spinless chain}

\section{Itinerant and single-ion magnetism}
In a uniform collinear Hartree state, an up electron feels $U n_\downarrow$.
Linearizing the spin-density difference gives $m=U\rho(E_F)m$ at zero
temperature for a smooth one-spin density of states. A finite cluster
needs a thermal derivative rather than a continuum density at a
single energy. Passing this uniform instability says nothing about
the preferred ordering wave vector.
\casestudy{stoner}{A fixed-density susceptibility test}
For a magnetic ion with $s\ge1$, $D(S^z)^2$ can split a multiplet even
in zero field. A term $E[(S^x)^2-(S^y)^2]$ mixes states whose $m_z$
differs by two. Stevens operators preserve the connection between
crystal symmetry and allowed tensor components. For a bond lacking
inversion symmetry the antisymmetric exchange
$\mathbf D_{ij}\cdot(\vS_i\times\vS_j)$ is allowed;
$\mathbf D_{ji}=-\mathbf D_{ij}$ is required by reversing the bond.
It is a two-site coupling, not a single-ion anisotropy.
\casestudy{anisotropy}{Zero-field splitting of a spin one}
\begin{exercise}[A frustrated triangle]\label{ex:triangle}
Find the energies and degeneracies of the Heisenberg triangle at
$J>0$. Reconstruct its zero-field partition function; the entropy
at $T\to0$ counts the full manifold, not just highest-weight states.
\end{exercise}
